\documentclass[11pt]{article}
\makeatletter\def\input@path{{../estilo/}}\makeatother
\usepackage{didistribuidores}
\renewcommand{\documentocorto}{Bitácora · Etapa 1}

\begin{document}
\portada{Bitácora de pruebas}{Etapa 1: monolito modular}%
  {Una aplicación y una base de datos, organizadas en módulos, con el flujo completo del pedido}%
  {Rama \cod{01-monolito}}{Octubre de 2026}

\tableofcontents
\newpage

\section{Qué se construyó en esta etapa}

En la etapa 1 construimos todo el sistema como \textbf{una sola aplicación con una sola base de
datos}, pero ordenada en módulos: \cod{usuarios}, \cod{restaurantes}, \cod{pedidos}, \cod{pagos} y
\cod{repartidores}. Cada módulo se divide en tres capas:

\begin{itemize}
  \item \cod{router.py}: recibe la petición HTTP y devuelve la respuesta;
  \item \cod{servicio.py}: las reglas del negocio;
  \item \cod{repositorio.py}: el SQL.
\end{itemize}

Un módulo solo habla con otro a través de su servicio, nunca leyendo sus tablas. Así, cuando en las
etapas 2 y 3 un módulo se convierte en servicio aparte, basta con cambiar una llamada a función por
una llamada de red; la lógica no se reescribe.

\begin{nota}[La decisión que atraviesa todas las pruebas]
Las garantías del sistema no descansan en el código de Python sino en la \textbf{base de datos}.
La idempotencia se apoya en una restricción \cod{UNIQUE}, y cada cambio de estado del pedido es un
\cod{UPDATE} que incluye el estado de origen en su condición. Cuando varias peticiones llegan al
mismo tiempo, la base decide quién gana y las demás reciben una respuesta clara.
\end{nota}

El pedido recorre los estados
\cod{PENDIENTE\_PAGO} $\rightarrow$ \cod{PAGADO} $\rightarrow$ \cod{ASIGNADO} $\rightarrow$
\cod{ENTREGADO}, o termina en \cod{CANCELADO} si no avanza a tiempo. Cada cambio sube en uno el
número de \emph{versión} del pedido, lo que en las pruebas sirve como contador: un pedido entregado
exactamente una vez termina con versión 4.

\begin{center}
\begin{tikzpicture}[node distance=1.6cm, every node/.style={font=\sffamily\small},
  est/.style={draw=azulsd, fill=grisfondo, rounded corners=3pt, minimum width=2.6cm, minimum height=0.8cm},
  fl/.style={-Stealth, thick, azulsd}]
  \node[est] (p) {PENDIENTE\_PAGO};
  \node[est, right=1.1cm of p] (pa) {PAGADO};
  \node[est, right=1.1cm of pa] (as) {ASIGNADO};
  \node[est, right=1.1cm of as] (en) {ENTREGADO};
  \node[est, below=1.2cm of pa, draw=rojo] (ca) {CANCELADO};
  \draw[fl] (p) -- node[above,font=\sffamily\scriptsize]{pago} (pa);
  \draw[fl] (pa) -- node[above,font=\sffamily\scriptsize]{acepta} (as);
  \draw[fl] (as) -- node[above,font=\sffamily\scriptsize]{código} (en);
  \draw[fl, rojo] (p) |- node[pos=0.25,left,font=\sffamily\scriptsize]{sin pago} (ca);
  \draw[fl, rojo] (pa) -- node[right,font=\sffamily\scriptsize]{sin repartidor (reembolso)} (ca);
\end{tikzpicture}
\end{center}

\subsection{Criterios del PRD que se comprueban}

\begin{tabularx}{\linewidth}{@{}l X l@{}}
\toprule
\textsf{\textbf{Historia}} & \textsf{\textbf{Criterio}} & \textsf{\textbf{Sección}} \\
\midrule
5.1 CA1 & Reintento con la misma \cod{Idempotency-Key}: mismo pedido, sin segundo cobro & \ref{sec:idem} \\
5.1 CA3 & La entrega se registra una sola vez aunque se confirme dos veces & \ref{sec:entrega} \\
5.3 CA1 & El repartidor solo ve pedidos \cod{PAGADO} de su zona & \ref{sec:aceptar} \\
5.3 CA2 & Dos repartidores aceptan a la vez: uno gana, el otro recibe un rechazo claro & \ref{sec:aceptar} \\
5.4 CA1--CA2 & Pedido pagado sin repartidor: se cancela y se reembolsa una sola vez & \ref{sec:timeouts} \\
\bottomrule
\end{tabularx}

\section{Entorno de pruebas}

Las pruebas se hicieron contra el sistema levantado con Docker Compose (PostgreSQL 17 y la
aplicación FastAPI en \cod{http://localhost:8000}). Para acortar los comandos se definieron unas
variables y funciones auxiliares:

\begin{comando}
B=http://localhost:8000
U=<id del usuario de prueba>      R=<id del restaurante>
P1=<id de la gringa, $60>         P2=<id del taco al pastor, $25>
nuevo_pedido() {   # crea un pedido con una clave nueva y devuelve su id
  curl -s -X POST $B/api/pedidos -H "Content-Type: application/json" \
    -H "Idempotency-Key: k-$(date +%s%N)" \
    -d "{\"usuario_id\":\"$U\",\"restaurante_id\":\"$R\",\"platillos\":[{\"platillo_id\":\"$P1\",\"cantidad\":1}]}" \
    | python3 -c "import sys,json; print(json.load(sys.stdin)['id'])"
}
psql_db() { docker compose exec -T db psql -U postgres -d didistribuidores -c "$1"; }
\end{comando}

\begin{nota}
Los scripts de la carpeta \cod{db/} (esquema, datos semilla y \cod{outbox}) solo se ejecutan cuando el
volumen de PostgreSQL está vacío. Para repetir las pruebas desde cero hay que usar
\cod{docker compose down -v}.
\end{nota}

\section{Arranque de la aplicación}

\queria Que la aplicación levanta, abre su conexión a la base de datos y responde.

\begin{comando}
docker compose up --build
curl http://localhost:8000/salud
\end{comando}

\esperado \cod{\{"estado":"ok"\}}.

\observado La respuesta fue \cod{\{"estado":"ok"\}}. Ese endpoint ejecuta un \cod{SELECT 1} contra la
base, así que una respuesta correcta confirma que la aplicación y la base están conectadas.

\figura[0.8]{img/fig-01.png}{Respuesta del endpoint \cod{/salud} con la aplicación levantada.}

\section{Usuarios y restaurantes}

\queria Que los módulos más sencillos funcionan con sus tres capas, y tener datos para probar los pedidos.

Como los datos semilla todavía no estaban listos, se insertó a mano, con \cod{psql}, un restaurante
(Tacos El Güero, zona 1) con dos platillos: el taco al pastor de \$25 y la gringa de \$60. Después
se consultó la lista de restaurantes, se creó un usuario (zona 1) y se consultaron el menú y el
usuario recién creado.

\observado La lista devolvió el restaurante insertado; el alta de usuario respondió con su
identificador y una fecha de creación en UTC; el menú mostró la gringa a 60.0 y el taco a 25.0.

En la misma captura aparece un detalle útil: en un intento se escribió el identificador del usuario
incompleto y el sistema respondió con un error de validación (\emph{no es un UUID válido}). Los
identificadores mal formados se rechazan antes de llegar a la base.

\figura{img/fig-02.png}{Alta del restaurante de prueba, consulta de restaurantes, creación del usuario, menú y validación de un identificador mal escrito.}

\section{Pedido idempotente}\label{sec:idem}

\queria El criterio 5.1 CA1: si la aplicación repite la petición con la misma
\cod{Idempotency-Key}, debe recibir el mismo pedido, sin que se cree otro ni se cobre de nuevo,
incluso cuando las peticiones llegan al mismo tiempo.

\comofunciona Cada petición de pedido lleva una clave de idempotencia en un encabezado. En la tabla de
pedidos hay una restricción \cod{UNIQUE (usuario\_id, clave\_idempotencia)} y el \cod{INSERT} usa
\cod{ON CONFLICT DO NOTHING}. Si llegan dos peticiones con la misma clave a la vez, la segunda espera a
que la primera termine y recibe el pedido ya creado. Además, el \textbf{total lo calcula el servidor}
con los precios del menú, nunca el cliente, y el pedido guarda el precio de cada platillo en el
momento en que se hizo.

\subsection{Prueba 1: un pedido y su reintento}

\begin{comando}
curl -i -X POST $B/api/pedidos -H "Content-Type: application/json" -H "Idempotency-Key: clave-1" \
  -d '{"usuario_id":"$U","restaurante_id":"$R","platillos":[
        {"platillo_id":"$P1","cantidad":2},{"platillo_id":"$P2","cantidad":1}]}'
\end{comando}

\esperado \cod{201 Created} la primera vez, con total 145 (2 × 60 + 25); \cod{200 OK} con el mismo
\cod{id} al repetir.

\observado La primera respuesta fue \cod{201 Created}, con el pedido en \cod{PENDIENTE\_PAGO}, un total
de 145.0, versión 1 y un código de entrega de seis dígitos. Al repetir exactamente la misma
petición, la respuesta fue \cod{200 OK} con el mismo identificador y el mismo contenido.

\figura{img/fig-03.png}{Primera petición: \cod{201 Created}, total 145.0 y estado \cod{PENDIENTE\_PAGO}.}
\figura{img/fig-04.png}{La misma petición con la misma clave: \cod{200 OK} y el mismo pedido.}

\subsection{Prueba 2: cinco peticiones al mismo tiempo}

\begin{comando}
for i in 1 2 3 4 5; do
  curl -s -o /dev/null -w "%{http_code}\n" -X POST $B/api/pedidos \
    -H "Content-Type: application/json" -H "Idempotency-Key: clave-2" -d "..." &
done; wait
psql_db "SELECT count(*) FROM pedidos.pedidos WHERE clave_idempotencia='clave-2'"
\end{comando}

\esperado Un \cod{201}, cuatro \cod{200} y \cod{count = 1}.

\observado Exactamente eso. Aunque las cinco peticiones llegaron juntas, solo se creó un pedido. Es
la prueba más importante de la idempotencia, porque demuestra que no hay un hueco entre
\emph{revisar si existe} y \emph{guardar}.

\figura{img/fig-05.png}{Cinco peticiones simultáneas con la misma clave: un 201, cuatro 200 y un solo pedido en la base.}

\subsection{Prueba 3: peticiones inválidas}

\begin{tabularx}{\linewidth}{@{}X l X@{}}
\toprule
\textsf{\textbf{Caso}} & \textsf{\textbf{Código}} & \textsf{\textbf{Mensaje}} \\
\midrule
Sin encabezado \cod{Idempotency-Key} & 400 & Falta el header Idempotency-Key \\
Platillo inexistente & 422 & Algún platillo no existe, no es de ese restaurante o no está disponible \\
Usuario inexistente & 404 & Usuario no encontrado \\
\bottomrule
\end{tabularx}

\figura[0.85]{img/fig-06.png}{Petición sin \cod{Idempotency-Key}: 400.}
\figura[0.85]{img/fig-07.png}{Platillo inexistente: 422.}
\figura[0.85]{img/fig-08.png}{Usuario inexistente: 404.}

\section{Pago}

\queria Que un pedido pasa de \cod{PENDIENTE\_PAGO} a \cod{PAGADO} y registra su cobro simulado una sola
vez, aunque la petición se repita o llegue varias veces al mismo tiempo.

\comofunciona Cambiar el estado y registrar el cobro ocurren en \textbf{una sola transacción}. El cambio
es un \cod{UPDATE ... WHERE estado = 'PENDIENTE\_PAGO'}: con peticiones simultáneas solo una consigue
el cambio; las demás encuentran que el pedido ya no está pendiente y reciben el mismo resultado.
Como segunda red de seguridad, la tabla de pagos tiene \cod{UNIQUE} sobre el pedido. Es la razón
principal por la que \cod{pagos} vive dentro del núcleo: el cobro y el cambio de estado nunca pueden
quedar a medias.

\subsection{Pago normal y repetido}
\begin{comando}
ID=$(nuevo_pedido)
curl -s -X POST $B/api/pedidos/$ID/pago
curl -s -X POST $B/api/pedidos/$ID/pago
curl -s $B/api/pagos/$ID
\end{comando}

\observado Los dos pagos respondieron 200 con el pedido en \cod{PAGADO} y la misma versión (2): el
segundo no cambió nada. La consulta del pago mostró un único cobro de 60.0 en estado \cod{COBRADO}.

\figura{img/fig-09.png}{Pago repetido del mismo pedido y consulta del cobro.}

\subsection{Cinco pagos al mismo tiempo}
\begin{comando}
ID=$(nuevo_pedido)
for i in 1 2 3 4 5; do curl -s -o /dev/null -w "%{http_code}\n" -X POST $B/api/pedidos/$ID/pago & done; wait
psql_db "SELECT (SELECT count(*) FROM pagos.pagos WHERE pedido_id='$ID') AS cobros,
                (SELECT version FROM pedidos.pedidos WHERE id='$ID') AS version"
\end{comando}

\observado Cinco respuestas 200, \cod{cobros = 1} y \cod{version = 2}: hubo un solo cobro y el pedido
cambió de estado una sola vez.

\figura{img/fig-10.png}{Cinco pagos simultáneos: cinco 200, un solo cobro y versión 2.}

\subsection{Pedido inexistente}
Pagar un pedido que no existe respondió \cod{404} con el mensaje \emph{Pedido no encontrado}.
\figura{img/fig-11.png}{Pago de un pedido inexistente: 404.}

\section{Aceptación del pedido por un repartidor}\label{sec:aceptar}

\queria Las historias 5.3 CA1 y CA2: el repartidor ve los pedidos pagados de su zona y acepta uno; si
dos aceptan al mismo tiempo, solo uno se queda con el pedido.

\comofunciona La asignación es un
\cod{UPDATE ... WHERE estado = 'PAGADO' AND repartidor\_id IS NULL}. Al primero que llega le cambia
la fila; al segundo ya no se le cumple la condición y recibe un conflicto. La regla de quién gana
vive en \cod{pedidos}, que es el dueño del estado del pedido; \cod{repartidores} solo guarda quién es
cada repartidor y en qué zona trabaja.

\subsection{Repartidores de prueba}
Se crearon tres: Ana y Luis en la zona 1, y Marta en la zona 2. Los tres respondieron \cod{201} con
\cod{disponible: true}.
\figura{img/fig-12.png}{Alta de los tres repartidores de prueba.}

\subsection{Pedidos disponibles por zona}
Se creó y pagó un pedido de la zona 1 y se consultaron los disponibles para Ana y para Marta.

\observado Ana recibió los pedidos pagados de su zona sin repartidor; Marta, de la zona 2, recibió
una lista vacía (\cod{[]}).
\figura{img/fig-13.png}{Pedidos disponibles para Ana (zona 1) y para Marta (zona 2).}

\subsection{Dos repartidores aceptan a la vez}
\begin{comando}
( curl -s -o /dev/null -w "Ana: %{http_code}\n"  -X POST $B/api/pedidos/$ID/aceptar \
    -H "Content-Type: application/json" -d "{\"repartidor_id\":\"$REP1\"}" &
  curl -s -o /dev/null -w "Luis: %{http_code}\n" -X POST $B/api/pedidos/$ID/aceptar \
    -H "Content-Type: application/json" -d "{\"repartidor_id\":\"$REP2\"}" &
  wait )
psql_db "SELECT estado, repartidor_id, version FROM pedidos.pedidos WHERE id='$ID'"
\end{comando}

\observado Ana recibió \cod{200} y Luis \cod{409}. El pedido quedó \cod{ASIGNADO}, con Ana como
repartidora y versión 3 (creado, pagado, asignado).
\figura{img/fig-14.png}{Aceptación simultánea: Ana 200, Luis 409 y el pedido \cod{ASIGNADO} con versión 3.}

\subsection{Otros casos}
Cuando Marta intentó aceptar un pedido de la zona 1 la respuesta fue \cod{422} (otra zona); aceptar un
pedido que aún no se había pagado respondió \cod{409}.
\figura{img/fig-15.png}{Aceptación desde otra zona (422) y de un pedido sin pagar (409).}

\section{Confirmación de la entrega}\label{sec:entrega}

\queria El criterio 5.1 CA3: con el código de entrega correcto, el pedido pasa a \cod{ENTREGADO} una sola
vez, aunque el repartidor lo envíe dos veces.

\comofunciona El código se comprueba dentro del mismo \cod{UPDATE} que cambia el estado:
\cod{WHERE estado = 'ASIGNADO' AND codigo\_entrega = ...}. Si el pedido ya estaba entregado con ese
mismo código, el servicio devuelve el mismo resultado sin modificar nada; la versión no sube.

\subsection{Preparar un pedido asignado}
Se creó un pedido, se pagó, lo aceptó un repartidor y se leyó su código (en esta prueba, 965343).
\figura{img/fig-16.png}{Pedido pagado, aceptado y lectura de su código de entrega.}

\subsection{Código incorrecto, correcto y repetido}
\observado Código incorrecto: \cod{422}. Correcto: \cod{200}. Repetido: \cod{200}, sin cambios.
\figura{img/fig-17.png}{Código incorrecto (422), correcto (200) y repetido (200).}

\subsection{Cinco entregas al mismo tiempo}
\observado Las cinco peticiones respondieron 200; el pedido quedó \cod{ENTREGADO} con
\textbf{versión 4}. Ese número prueba que la entrega ocurrió una sola vez: el pedido pasó exactamente
por cuatro versiones. Si se hubiera entregado más de una vez, la versión sería mayor.
\figura{img/fig-18.png}{Cinco entregas simultáneas: cinco 200, \cod{ENTREGADO} y versión 4.}

\subsection{Entregar un pedido sin repartidor}
Confirmar la entrega de un pedido pagado pero sin repartidor respondió \cod{409}.
\figura{img/fig-19.png}{Entrega de un pedido sin repartidor: 409.}

\section{Cancelación por tiempo y reembolso}\label{sec:timeouts}

\queria La historia 5.4: un pedido que nadie entrega se cancela y se reembolsa, una sola vez, y nunca
queda en el limbo.

\comofunciona Una tarea en segundo plano, que arranca junto con la aplicación, revisa cada 5 segundos
los pedidos vencidos. Hay \textbf{dos relojes}:

\begin{itemize}
  \item \cod{TIMEOUT\_PAGO\_SEG} (30 s en la demo, 10 min normal): un pedido sin pagar pasa a
        \cod{CANCELADO} con motivo \cod{sin\_pago}, sin reembolso porque nunca hubo cobro.
  \item \cod{TIMEOUT\_ASIGNACION\_SEG} (60 s en la demo, 20 min normal): un pedido pagado sin
        repartidor pasa a \cod{CANCELADO} con motivo \cod{sin\_repartidor} y su pago queda
        \cod{REEMBOLSADO}.
\end{itemize}

Cancelar y reembolsar son \textbf{una sola transacción}; y como el \cod{UPDATE} exige el estado de
origen, un pedido ya cancelado no vuelve a procesarse: nunca se reembolsa dos veces. Los pedidos que
ya tienen repartidor no se cancelan.

\subsection{Pedido que nunca se paga}
Se creó un pedido y se esperaron 40 segundos.

\observado \cod{CANCELADO} con motivo \cod{sin\_pago}; en el registro aparece
\emph{Timeouts: 1 pedido(s) cancelado(s)}.
\figura{img/fig-20.png}{Pedido sin pagar: \cod{CANCELADO} con motivo \cod{sin\_pago}.}

\subsection{Pedido pagado que nadie acepta}
Se creó un pedido, se pagó y se esperaron 70 segundos.

\observado \cod{CANCELADO} con motivo \cod{sin\_repartidor} y el pago \cod{REEMBOLSADO}. Las fechas del
propio pago lo confirman: cobrado a las 02:38:03 y reembolsado a las 02:39:06, unos 63 segundos
después, coherente con el reloj de 60 segundos más la revisión cada 5.
\figura{img/fig-21.png}{Pedido pagado sin repartidor: \cod{CANCELADO} y pago \cod{REEMBOLSADO}.}

\subsection{Un pedido asignado no se cancela}
Se creó un pedido, se pagó, lo aceptó un repartidor y se esperaron 75 segundos: siguió \cod{ASIGNADO}.
\figura{img/fig-22.png}{Pedido asignado después de 75 segundos: sigue \cod{ASIGNADO}.}

\subsection{Pagar un pedido ya cancelado}
Respondió \cod{409}. Si el reloj y el cliente compiten, gana quien llega primero a la base de datos.
\figura[0.85]{img/fig-23.png}{Pago de un pedido cancelado: 409.}

\section{Infraestructura de la etapa 1}

\queria Que el entorno se levanta desde cero con un solo comando, con datos de prueba y con la tabla
donde se guardan los eventos.

El \cod{docker-compose.yml} usa PostgreSQL 17, un volumen con nombre y un \emph{healthcheck} en la
base; la aplicación declara \cod{depends\_on} con \cod{condition: service\_healthy}, así que arranca
hasta que la base está realmente lista.

\subsection{Levantar desde cero}
\begin{comando}
docker compose down -v
docker compose up --build -d
docker compose ps
\end{comando}
\figura{img/fig-24.png}{Sistema levantado desde cero: base \cod{healthy} y aplicación \cod{Up}.}

\subsection{Datos semilla}
3 restaurantes, 9 platillos y 5 repartidores, tal como define \cod{db/02-semilla.sql}.
\figura{img/fig-25.png}{Datos semilla cargados.}

\subsection{Tabla \cod{outbox}}
\cod{pedidos.outbox} tiene las columnas \cod{id}, \cod{tipo}, \cod{payload}, \cod{creado\_en} y
\cod{publicado\_en}, y un índice parcial, \cod{outbox\_pendientes\_idx}, que solo cubre las filas con
\cod{publicado\_en} nulo. Ese índice permite al publicador de la etapa 2 encontrar rápido los
eventos pendientes.
\figura{img/fig-26.png}{Estructura de la tabla \cod{pedidos.outbox} y sus índices.}

\subsection{Eventos guardados por un pedido completo}
\observado Tres eventos, en orden: \cod{pedido.pagado} (versión 2), \cod{pedido.asignado} (versión 3)
y \cod{pedido.entregado} (versión 4). Cada evento se guarda en la misma transacción que el cambio de
estado: no puede existir un cambio sin su evento ni un evento sin su cambio. En esta etapa quedan con
\cod{publicado\_en} nulo, porque el publicador hacia RabbitMQ llega en la etapa 2.
\figura{img/fig-27.png}{Tres eventos en el \cod{outbox} para un pedido completo, con versiones 2, 3 y 4.}

\section{Resumen de la etapa 1}

\begin{longtable}{@{}p{4.6cm} p{6.6cm} l l@{}}
\toprule
\textsf{\textbf{Qué se comprobó}} & \textsf{\textbf{Resultado observado}} & \textsf{\textbf{Fig.}} & \textsf{\textbf{Estado}} \\
\midrule
La aplicación responde y está conectada a la base & \cod{/salud} devuelve ok & 1 & \ok \\
Usuarios, restaurantes y menús & Se crean y consultan; los ids mal formados se rechazan & 2 & \ok \\
Reintento con la misma clave & 201 y luego 200 con el mismo id & 3--4 & \ok \\
Cinco peticiones simultáneas con la misma clave & Un 201, cuatro 200 y un solo pedido & 5 & \ok \\
Peticiones inválidas & 400, 422 y 404 con mensajes claros & 6--8 & \ok \\
Pago repetido y simultáneo & Un solo cobro y versión 2 & 9--11 & \ok \\
Dos repartidores aceptan a la vez & Un 200 y un 409; asignado con versión 3 & 12--15 & \ok \\
Entrega repetida y simultánea & Una sola entrega, versión 4 & 16--19 & \ok \\
Cancelación por tiempo & \cod{sin\_pago} sin reembolso; \cod{sin\_repartidor} con reembolso; el asignado no se cancela & 20--23 & \ok \\
Infraestructura y \cod{outbox} & Base sana, datos semilla, tabla \cod{outbox} y tres eventos & 24--27 & \ok \\
\bottomrule
\end{longtable}

\subsection*{Lo que dejamos preparado para las siguientes etapas}
\begin{itemize}
  \item \cod{pedidos.repartidor\_id} es un identificador \textbf{sin llave foránea}, para que separar
        \cod{repartidores} en la etapa 3 no obligue a romper restricciones.
  \item Cada cambio de estado ya deja su evento en \cod{pedidos.outbox}: la etapa 2 solo tiene que
        publicarlos.
  \item Los relojes y el intervalo de revisión se configuran por variables de entorno, para acortarlos
        en la demo.
\end{itemize}

\section{Quién hizo qué}

El trabajo se repartió desde el inicio en dos roles: \textbf{Erick} (Persona A) con las etapas 1 y 3, y
\textbf{Yuki} (Persona B) con las etapas 2 y 4. Cada cambio lo revisó el otro antes de entrar a la rama
de entrega.

\begin{tabularx}{\linewidth}{@{}p{4.2cm} X@{}}
\toprule
\textsf{\textbf{Integrante}} & \textsf{\textbf{Aportes en la etapa 1}} \\
\midrule
Erick Isaac Echeverria Goicochea (@ErickI07) & Estructura de módulos y esquema SQL del núcleo; arranque de la aplicación, pool de conexiones y manejo de errores; módulos de usuarios y restaurantes; creación de pedidos con idempotencia; pagos; repartidores y aceptación de pedidos; confirmación de la entrega; cancelación por tiempo con reembolso; pruebas de los pasos 1 a 7; pull request de la etapa. \\
Yukioayax Canek Gabriel Hernández (@YukiGab) & Docker Compose final (PostgreSQL 17, volúmenes y \emph{healthchecks}), datos semilla y tabla \cod{outbox}; pruebas de la infraestructura; revisión del pull request. \\
Los dos & Decisiones de alcance y de diseño del PRD; contrato de eventos; revisión cruzada; evidencias y bitácora. \\
\bottomrule
\end{tabularx}

\end{document}
